\documentclass[12pt,a4paper]{article}

% ABNT formatting packages
\usepackage[inner=30mm,outer=20mm,top=30mm,bottom=20mm]{geometry}
\usepackage[utf8]{inputenc}
\usepackage[T1]{fontenc}
\usepackage[brazilian]{babel}
\usepackage{graphicx}
\usepackage{float}
\usepackage[colorlinks=true,linkcolor=black,urlcolor=black,citecolor=black]{hyperref}
\usepackage{indentfirst}
\usepackage{setspace}
\usepackage{helvet}
\usepackage[font=normalsize,labelfont=bf,labelsep=endash,justification=centering,position=above]{caption}

\renewcommand{\familydefault}{\sfdefault}
\onehalfspacing
\setlength{\parindent}{1.5cm}

\input{../../../../out/drive-links.tex}

\begin{document}

\section{Diagramas de Sequência -- Usuário}

Nas figuras a seguir são apresentados os diagramas de sequência das operações relacionadas
ao domínio de usuários e autenticação, descrevendo o fluxo de mensagens entre os componentes do sistema.

\subsection{Usuário}

\begin{figure}[H]
    \caption{Diagrama de sequência -- Criar Usuário}
    \label{fig:seq-create-user}
    \centering
    \includegraphics[width=0.95\textwidth,height=0.75\textheight,keepaspectratio]{../../../../out/src/plantuml/sequencia/user/create_user.png}

    {\footnotesize \textbf{Fonte:} O autor.
    \textbf{Acesso:} \href{\linkCreateUser}{Diagrama de Sequência: Criar Usuário}.}
\end{figure}

\begin{figure}[H]
    \caption{Diagrama de sequência -- Buscar Usuário}
    \label{fig:seq-get-user}
    \centering
    \includegraphics[width=0.95\textwidth,height=0.75\textheight,keepaspectratio]{../../../../out/src/plantuml/sequencia/user/get_user.png}

    {\footnotesize \textbf{Fonte:} O autor.
    \textbf{Acesso:} \href{\linkGetUser}{Diagrama de Sequência: Buscar Usuário}.}
\end{figure}

\begin{figure}[H]
    \caption{Diagrama de sequência -- Atualizar Usuário}
    \label{fig:seq-update-user}
    \centering
    \includegraphics[width=0.95\textwidth,height=0.75\textheight,keepaspectratio]{../../../../out/src/plantuml/sequencia/user/update_user.png}

    {\footnotesize \textbf{Fonte:} O autor.
    \textbf{Acesso:} \href{\linkUpdateUser}{Diagrama de Sequência: Atualizar Usuário}.}
\end{figure}

\begin{figure}[H]
    \caption{Diagrama de sequência -- Atualizar Status do Usuário}
    \label{fig:seq-update-user-status}
    \centering
    \includegraphics[width=0.95\textwidth,height=0.75\textheight,keepaspectratio]{../../../../out/src/plantuml/sequencia/user/update_user_status.png}

    {\footnotesize \textbf{Fonte:} O autor.
    \textbf{Acesso:} \href{\linkUpdateUserStatus}{Diagrama de Sequência: Atualizar Status do Usuário}.}
\end{figure}

\begin{figure}[H]
    \caption{Diagrama de sequência -- Excluir Usuário}
    \label{fig:seq-delete-user}
    \centering
    \includegraphics[width=0.95\textwidth,height=0.75\textheight,keepaspectratio]{../../../../out/src/plantuml/sequencia/user/delete_user.png}

    {\footnotesize \textbf{Fonte:} O autor.
    \textbf{Acesso:} \href{\linkDeleteUser}{Diagrama de Sequência: Excluir Usuário}.}
\end{figure}

\subsection{Autenticação}

\begin{figure}[H]
    \caption{Diagrama de sequência -- Login}
    \label{fig:seq-login}
    \centering
    \includegraphics[width=0.95\textwidth,height=0.75\textheight,keepaspectratio]{../../../../out/src/plantuml/sequencia/auth/login.png}

    {\footnotesize \textbf{Fonte:} O autor.
    \textbf{Acesso:} \href{\linkLogin}{Diagrama de Sequência: Login}.}
\end{figure}

\begin{figure}[H]
    \caption{Diagrama de sequência -- Filtro JWT}
    \label{fig:seq-jwt-filter}
    \centering
    \includegraphics[width=0.95\textwidth,height=0.75\textheight,keepaspectratio]{../../../../out/src/plantuml/sequencia/auth/jwt_filter.png}

    {\footnotesize \textbf{Fonte:} O autor.
    \textbf{Acesso:} \href{\linkJwtFilter}{Diagrama de Sequência: Filtro JWT}.}
\end{figure}

\end{document}
